\section{The determinant class and the canonical metric}\label{sec:canonical}

Write $R=2r_0$, $c=2aH$ and $h=c_1(H)=[\omega]$.
Let $\Omega$ be the K\"ahler form of the harmonic metric
on $M_{\reg}$. In the normalization of
Biswas--Schumacher, denote the Weil--Petersson form by
$\omega^{\WP}$. There is a fixed positive constant $b$
such that
\begin{equation}\label{eq:WP-normalization}
 \Omega=b\,\omega^{\WP}.
\end{equation}
This constant only converts trace and real-metric conventions;
the cone and actual metrics throughout the paper use the
same harmonic convention. We retain $b$ explicitly in
the class calculation.

\subsection{The positive determinant polarization}
Choose a sufficiently divisible positive integer $k$, a
divisor $D\in|H^k|$, and a point $x\in S$. For the fixed
Mukai vector $v=2v_0$ define
\begin{equation}\label{eq:determinant-class}
 u_D=-R[\OO_D]+\chi(F\otimes\OO_D)[\OO_x],
 \qquad
 \lambda_D=\det Rp_*(\mathcal E\otimes q_S^*u_D).
\end{equation}
The Euler characteristic depends only on $v$, not on $F$.
The expression on the right is the determinant of cohomology
on family parameter spaces, descended to $M$. It is the
positive convention of
\cite[Definition~8.1.9]{hl97}. Here $p$ and $q_S$ are the
projections to the parameter space and $S$. Any positive
power required in the construction can be absorbed into
$\lambda_D$, with the reciprocal change in the scalar
class coefficient below. The auxiliary $k$ does not change
the fixed form $\omega$ used for the harmonic metric.

The stabilizer weights also make descent explicit in this
family. For a stable factor $G$ of vector $v_0$, additivity
of the Euler characteristic gives
\[
 \chi(G\otimes u_D)
 =-2r_0\chi(G\otimes\OO_D)
             +2\chi(G\otimes\OO_D)r_0=0.
\]
Each multiplicity-group determinant character is therefore
trivial at a closed polystable orbit, including a repeated
summand. At a stable point the same equality follows directly
from the definition of $u_D$.

\begin{lemma}\label{lem:ample}
The determinant line $\lambda_D$ is ample on $M$.
\end{lemma}

\begin{proof}
The determinant construction is semiample and defines
the projective Donaldson--Uhlenbeck morphism
\cite[Definition and Theorem~8.2.8]{hl97}. Its closed-point
invariant is the reflexive hull of the slope graded object
together with the defect cycle of that graded object
\cite[Theorem~8.2.11]{hl97}. Proposition~\ref{prop:no-bubbling}
says that every slope graded factor here is locally free.
Thus the defect cycle is zero. The remaining invariant
is the polystable bundle itself, and slope S-equivalence
agrees with Gieseker S-equivalence. The morphism is
therefore injective on closed points. Being projective
and quasi-finite, it is finite. A positive power of
$\lambda_D$ is the pullback of an ample line under this
finite morphism, proving ampleness. This argument needs
neither normality of the target nor an isomorphism with it.
\end{proof}

\begin{lemma}\label{lem:class}
With the determinant convention \eqref{eq:determinant-class},
before replacing it by a descent power,
\begin{equation}\label{eq:omega-class}
 [\Omega]=\frac{4\pi^2b}{Rk}\,c_1(\lambda_D)
       \quad\hbox{in }H^2(M_{\reg},\RR).
\end{equation}
\end{lemma}

\begin{proof}
The Chern character of the $K$-theory class in
\eqref{eq:determinant-class} is
\begin{equation}\label{eq:chu}
 \ch(u_D)=(0,-Rkh,k(c\cdot h)[x]).
\end{equation}
Indeed, $\ch(\OO_D)=1-e^{-kh}$ and
$\chi(F\otimes\OO_D)=k(c\cdot h)-Rk^2h^2/2$
on a K3 surface.

Work first on a stable parameter space with a universal
family $\mathcal E$. Since $H^1(S,\RR)=0$,
\[
 c_1(\mathcal E)=q_S^*c+p^*\beta
\]
for a class $\beta$ on that space. Grothendieck--Riemann--Roch
and \eqref{eq:chu} give
\begin{align}
 c_1(\lambda_D)
 &=-Rk\,p_*(\ch_2(\mathcal E)\,q_S^*h)
                  +k(c\cdot h)\beta \notag\\
 &=-\frac{k}{2}\,p_*(\ch_2(\End\mathcal E)\,q_S^*h).
                                                    \label{eq:GRR}
\end{align}
For the second equality use
$\ch_2(\End\mathcal E)=2R\ch_2(\mathcal E)-c_1(\mathcal E)^2$
and
$p_*(c_1(\mathcal E)^2q_S^*h)=2(c\cdot h)\beta$.
This agrees with the discriminant formula in
\cite[Proposition~8.3.1(ii)]{hl97}.

The fixed-determinant curvature identity of
\cite[Proposition~4.3, equation~(4.6)]{biswas16}, in
surface dimension, is
\begin{equation}\label{eq:WP-curvature}
 \frac{\omega^{\WP}}{2\pi^2}
   =-\frac1R\,p_*(\ch_2(\End\mathcal E,h_E)\wedge\omega).
\end{equation}
The fixed-determinant hypothesis holds because
$\operatorname{Pic}^0(S)=0$. Comparing
\eqref{eq:GRR} and \eqref{eq:WP-curvature} gives the stated
coefficient.

One can check the sign through the Quillen line as well.
In the convention of \cite[Theorem~4.6, equations~(4.11)
and~(4.13)]{biswas16}, the positive line is the dual
determinant
\begin{equation}\label{eq:Quillen}
 Q_D=\left[\det Rp_*\left(
 (\End\mathcal E-R^2\OO)\otimes
                   (q_S^*H^k-q_S^*H^{-k})\right)\right]^{-1}.
\end{equation}
Its first Chern class is
\[
 c_1(Q_D)=-2k\,p_*(\ch_2(\End\mathcal E)\,q_S^*h)
                  =4c_1(\lambda_D).
\]
Using the nondual determinant would reverse this sign.

Finally these computations give equality on the coarse
stable locus, not merely on a family. The stable Quot
parameter space is a principal $\operatorname{PGL}(m)$
bundle over $M^s$ for an appropriate framing dimension
$m$. Its universal endomorphism bundle and determinant
classes descend. Pullback in real degree-two cohomology
is injective, by the low-degree Leray sequence and
$H^1(\operatorname{PGL}(m),\RR)=0$. Thus the equality
proved on that atlas gives \eqref{eq:omega-class} on
$M^s=M_{\reg}$. No universal bundle on the full coarse
moduli space is required.
\end{proof}

Put $\lambda=\lambda_D$ and
$\kappa=4\pi^2b/(Rk)$, adjusted if a positive descent
power has been taken. Choose a smooth K\"ahler
representative $\theta\in c_1(\lambda)$, obtained from
a projective embedding and scaled by its degree.

\subsection{Continuous potentials in the global class}
Theorem~11 of \cite{buchdahl21} applies to polystable
bundles deformable to stable bundles over a compact Hodge
manifold. The surface $S$ is Hodge since $[\omega]=c_1(H)$;
all our polystable objects are bundles; and the simple
points in Section~\ref{subsec:quivers} show local density
of stable bundles. The analytic moduli germs are those
identified in Section~\ref{sec:harmonic}. That theorem
therefore gives continuous plurisubharmonic local
potentials for $\Omega$. We now check that they define
a current in the global class from Lemma~\ref{lem:class}.

The variety $M$ is Cohen--Macaulay and
$\operatorname{codim}\Sing M\ge6$. Scheja's restriction
theorem \cite[Section~5, Corollary after Satz~III]{scheja61}
gives
\begin{equation}\label{eq:Scheja}
 H^1(M,\OO_M)\simeq H^1(M_{\reg},\OO_{M_{\reg}}).
\end{equation}
In the theorem's depth convention, for degree one it
suffices that
$\operatorname{depth}_x\OO_M\ge\dim_x\Sing M+3$;
our dimension and codimension bounds satisfy this.
Combining \eqref{eq:Scheja} with
\eqref{eq:forms-cohomology} gives vanishing of the
right side.

By Lemma~\ref{lem:class}, there is a smooth real
one-form $\alpha$ on $M_{\reg}$ with
$d\alpha=\Omega-\kappa\theta$. Its $(0,1)$ part is
$\dbar$-closed. The preceding vanishing and the
Dolbeault theorem give a smooth function $f$ with
$\alpha^{0,1}=\dbar f$. Since $\alpha$ is real,
\begin{equation}\label{eq:global-potential}
 \Omega=\kappa\theta+dd^c\varphi
 \quad\hbox{on }M_{\reg},\qquad
 \varphi=2\operatorname{Im}f.
\end{equation}

Let $U$ be a small neighbourhood in $M$. Write
$\Omega=dd^cu$ on $U_{\reg}$, where $u$ is the
continuous local potential supplied above, and
$\theta=dd^cw$ with $w$ a smooth local potential.
Then
\[
 h_U=\varphi-u+\kappa w
\]
is pluriharmonic on $U_{\reg}$. Both $h_U$ and
$-h_U$ are plurisubharmonic there. The
Grauert--Remmert extension theorem across a subset of
complex codimension at least two in a normal space
\cite[Section~1.7, Satz~4]{gr56} extends both as
plurisubharmonic functions on $U$. In this codimension
the theorem requires no prior upper bound. Their sum
is zero on the dense regular locus, hence everywhere
by uniqueness of plurisubharmonic extension. The
extensions are consequently finite and continuous:
each is both upper and lower semicontinuous. Their
$dd^c$ currents are opposite positive currents, so
both vanish.

It follows that $\varphi=u-\kappa w+h_U$ extends
continuously over $U$. The extensions agree on overlaps
by density, producing a global continuous function
$\varphi$ on $M$. Formula \eqref{eq:global-potential}
now defines a positive closed current on $M$ in
$\kappa c_1(\lambda)$, with the claimed local
continuous potentials.

\subsection{Identification by Ricci-flat uniqueness}
The harmonic metric on the stable locus is
hyperk\"ahler. This is Itoh's theorem for irreducible
Einstein--Hermitian connections on a K3 surface
\cite[Theorem~2.1 and Remark~(iii)]{itoh88}; the
Einstein--Hermitian extension in that remark allows
the nonzero slope here. The associated holomorphic
symplectic form, in harmonic coordinates, is a fixed
nonzero multiple of
\begin{equation}\label{eq:Mukai-form}
 (a,b)\longmapsto\int_S\tr(a\wedge b)\wedge\sigma.
\end{equation}
Indeed, the quaternionic complex structures on
endomorphism-valued one-forms and the trace $L^2$
pairing give \eqref{eq:Mukai-form} by contraction
with the parallel form $\sigma$. This is the
Serre-duality form defining the Mukai symplectic
structure, so it agrees with $\sigma_M$ up to the
fixed normalization. This identification is a
calculation of the forms; it does not require
extending a Riemannian tensor across $\Sigma$.

Set $d=\dim_{\CC}M=4q+2$. The hyperk\"ahler volume
identity on $M_{\reg}$ says that $\Omega^d$ is a
constant positive multiple of the adapted measure
defined by $\sigma_M^{d/2}$ and its conjugate. The
constant is determined by the total volume of
$\kappa c_1(\lambda)$. The current has locally
bounded potentials, so its Monge--Amp\`ere measure
does not charge the singular analytic set. The
volume identity therefore holds on all of $M$.

The variety is projective, has canonical
singularities and trivial canonical bundle. The
class $\kappa c_1(\lambda)$ is K\"ahler and belongs
to the real N\'eron--Severi space. The existence and
uniqueness theorem of Eyssidieux--Guedj--Zeriahi
\cite[Theorem~7.5 and Corollary~7.6]{egz09} applies
to precisely these data. It identifies our current
with the canonical singular Ricci-flat current in
that class. This proves part~(i) of
Theorem~\ref{thm:main}. Together with
Section~\ref{sec:distances}, it proves all claims
of Corollary~\ref{cor:completion-main} for the
canonical metric.
